% Ad Atticum 7.16 (ad-atticum-07-16)
% Source: https://raw.githubusercontent.com/PerseusDL/canonical-latinLit/master/data/phi0474/phi057/phi0474.phi057.perseus-lat1.xml
% Fetched by scripts/fetch_latin.py

% Dateline (Perseus): Scr. Calibus iv K. Febr. a. 705 (49).

\ciceroLetterOpener{CICERO ATTICO salutem}

\ciceroSection{1} omnis arbitror mihi tuas litteras redditas esse sed primas praepostere, reliquas ordine quo sunt missae per Terentiam. de mandatis Caesaris adventuque Labieni et responsis consulum ac Pompei scripsi ad te litteris iis quas a. d. v Kal. Capua dedi pluraque praeterea in eandem epistulam conleci.

\ciceroSection{2} nunc has exspectationes habemus duas, unam quid Caesar acturus sit cum acceperit ea quae referenda ad illum data sunt L. Caesari, alteram quid Pompeius agat. qui quidem ad me scribit paucis diebus se firmum exercitum habiturum spemque adfert, si in Picenum agrum ipse venerit, nos Romam redituros esse. Labienum secum habet non dubitantem de imbecillitate Caesaris copiarum; cuius adventu Gnaeus noster multo animi plus habet. nos a consulibus Capuam venire iussi sumus ad Nonas Febr. Capua profectus sum Formias a. d. iii Kal. eo die cum Calibus tuas litteras hora fere nona accepissem, has statim dedi. de Terentia et Tullia tibi adsentior.

\ciceroSection{3} ad quas scripseram ad te ut referrent. si nondum profectae sunt, nihil est quod se moveant, quoad perspiciamus quo loci sit res.
